% !PRINTED 381
\weberpage{381}

%% [running head, not body text -- commented out 2026-08-07] Page furniture (running
%% head, folio, signature) is excluded from this edition by standing policy.
%% \noindent\S.\,123. \hfill \textsc{Reducirte Zahlen.} \hfill 381

%<<CONT
\weq{6}{\omega = \alpha - \dfrac{1}{\omega_1},}

so dass auch $\omega_1 = \xi_1 + i\eta_1$ mit $\omega$ äquivalent ist. Die Zahl $\omega_1$
behandeln wir nun wieder ebenso wie $\omega$, indem wir, wenn $\omega_1 - \alpha_1$
noch nicht reducirt ist,

\weq{7}{\omega_1 = \alpha_1 - \dfrac{1}{\omega_2}}

setzen u.\,s.\,f. Betrachten wir die Reihe der nach Analogie von (5) gebildeten
Grössen
\[
r_1^2 = (\xi_1-\alpha_1)^2 + \eta_1^2, \quad r_2^2 = (\xi_2-\alpha_2)^2 + \eta_2^2 \ldots,
\]
so kommt es also jetzt nur darauf an, nachzuweisen, dass wir nach einer
endlichen Zahl von Schritten dieser Art zu einem $r_\nu^2$ kommen, das gleich
oder grösser als 1 ist.

Nach (6) ist aber
\[
r^2 = (\omega-\alpha)(\omega'-\alpha) = \dfrac{1}{\xi_1^2+\eta_1^2},
\]
und danach ergiebt die Vergleichung der imaginären Theile auf beiden
Seiten von (6):
\[
\eta = r^2 \eta_1.
\]
Ebenso folgt die Reihe der Gleichungen

\weq{8}{\eta = r^2\eta_1, \quad \eta_1 = r_1^2\eta_2, \quad \eta_2 = r_2^2\eta_3 \ldots}

Nun ist aber $i\eta = y\sqrt{d}$, und also nach der Formel (6), \S.\,122
\[
i\eta = \dfrac{\sqrt{D}}{2c},
\]
worin $c$ eine ganze Zahl ist. Weil aber äquivalente Zahlen dieselbe
Discriminante haben, so folgt ebenso:
\[
i\eta_1 = \dfrac{\sqrt{D}}{2c_1}, \quad i\eta_2 = \dfrac{\sqrt{D}}{2c_2}, \ldots,
\]
worin $c, c_1, c_2 \ldots$ eine Reihe positiver ganzer Zahlen ist. Danach
erhält man aus (8):

\weq{9}{c_1 = r^2 c, \quad c_2 = r_1^2 c_1, \quad c_3 = r_2^2 c_2 \ldots}

So lange aber die Zahlen $r^2, r_1^2 \ldots r_\nu^2$ kleiner als 1 sind, folgt hieraus

\weq{10}{c > c_1 > c_2 \ldots > c_{\nu+1},}

und weil es nur eine endliche Anzahl positiver Zahlen geben kann, die
kleiner als $c$ sind, so muss diese Reihe abbrechen, womit unser Satz 1.
bewiesen ist.
